\chapter{Skew Brownian Motion}\label{ch:sbm}

Let $\beta\in[-1,1]$. A \emph{skew Brownian motion} of parameter $\beta$
started at $x$ is a solution of
\begin{equation}\label{eq:sbm}
	X_t = x + W_t + \beta\hat{L}^0_t(X),
\end{equation}
where $W$ is a Brownian motion and $\hat{L}^0(X)$ is the symmetric local time
of $X$ at $0$ (Definition~\ref{def:symlocaltime}). This process was introduced
by \citet{ItoMcKean1965} and studied by \citet{Walsh1978} and
\citet{HarrisonShepp1981}; a detailed survey is \citet{Lejay2006}. For
$\beta = 0$ it is a Brownian motion, and for $\beta = \pm1$ it is a reflected
Brownian motion.

When $|\beta|<1$, \eqref{eq:sbm} is the equation \eqref{eq37} of
Chapter~\ref{ch:local-times-unknown} with $\phi\equiv1$ and
$\nu = \beta\delta_0\in M(\mathbb{R})$, so by Theorem~\ref{th:singular} it has a
unique strong solution, and pathwise uniqueness holds; the same is true for
$|\beta|\leqslant1$ by \citet{HarrisonShepp1981}. It is also the equation of
Remark~\ref{rem:alpha-sbm} with a constant coefficient. In particular $X$ is a
strong Markov process. Throughout this chapter, $\mathbb{P}_x$ and $\E_x$
refer to the solution started at $x$, and we write $\hat{L}_t = \hat{L}^0_t(X)$.

\begin{remark}\label{rem:sbm-scale}
	For $\nu = \beta\delta_0$ with $|\beta|<1$, the function $f_\nu$ of
	Lemma~\ref{lem:fnu} is $f_\nu = 1$ on $]-\infty,0[$ and
	$f_\nu = \frac{1-\beta}{1+\beta}$ on $[0,\infty[$. By
	Proposition~\ref{prop31}, $F_\nu(X)$ is then a local martingale, where
	\[
		F_\nu(x) = x\,1_{\{x<0\}} + \frac{1-\beta}{1+\beta}\,x\,1_{\{x\geqslant0\}} .
	\]
	The slopes of this scale function satisfy
	$(1+\beta)F'_\nu(0+) = (1-\beta)F'_\nu(0-)$; this is the transmission
	condition that appears in Section~\ref{sec:sbm-pde}.
\end{remark}

In this chapter we follow \citet{Lejay2006} and \citet{EtoreMartinez2012}. We
compute the joint law of $X_t$ and $\hat{L}_t$, deduce the transition density
of $X$, and describe the partial differential equation it solves.

\section{The law of skew Brownian motion}

\subsection{Construction by random flipping of excursions}
We first describe a construction of skew Brownian motion that shows its
relation to Brownian motion \citep[see][Chapter~XII, Exercise~2.16]{RevuzYor1999}.
Starting from a reflected Brownian motion, we change the sign of each excursion
independently, keeping it positive with probability $\frac{1+\beta}{2}$.

Let $B$ be a Brownian motion started at $0$, and let $(Y_n)_{n\geqslant0}$ be
a sequence of independent random variables, independent of $B$, with
\[
	\mathbb{P}(Y_n = 1) = \frac{1+\beta}{2}, \qquad \mathbb{P}(Y_n = -1) = \frac{1-\beta}{2}.
\]
The set $\{t : B_t\neq0\}$ is a countable union of disjoint open intervals, the
excursion intervals of $B$. They can be enumerated measurably as
$(I_n)_{n\geqslant0}$: for instance, for each $k$ and $j$ only finitely many
of them are contained in $[0,k]$ and have length at least $2^{-j}$, and these
can be listed by their left end-points. Define
\[
	X_t = Y_n\,|B_t| \quad\text{if } t\in I_n, \qquad X_t = 0 \quad\text{if } B_t = 0 .
\]
Thus $|X| = |B|$, and each excursion of $X$ is an excursion of $|B|$ given a
random sign. The process $X$ is a skew Brownian motion of parameter $\beta$
started at $0$: there is a Brownian motion $W$ such that $X$ satisfies
\eqref{eq:sbm} with $x = 0$ \citep{HarrisonShepp1981,Lejay2006}.

The symmetric local time of $X$ is a functional of $|X|$, since
$\hat{L}_t = \lim_{\varepsilon\to0}\frac{1}{2\varepsilon}\int_0^t1_{\{|X_s|<\varepsilon\}}\,ds$
by the occupation formula. As $|X| = |B|$, we get
$\hat{L}_t = L^0_t(B)$.

\begin{remark}
	\citet{Lejay2006} describes extensions of this ``random flipping of
	excursions'' method, which also give a decomposition of the Itô measure of
	the excursions of $X$.
\end{remark}

\subsection{Joint law of skew Brownian motion and its local time}

We use the notation
\[
	p_t(z) = \frac{1}{\sqrt{2\pi t}}\exp\Bigl(-\frac{z^2}{2t}\Bigr),
	\qquad
	q(a,t) = \frac{a}{\sqrt{2\pi t^3}}\exp\Bigl(-\frac{a^2}{2t}\Bigr) = \frac{a}{t}\,p_t(a),
\]
for $t>0$, $z\in\mathbb{R}$ and $a>0$. Recall that $q(a,\cdot)$ is the density
of the hitting time $T_a$ of $a$ by a Brownian motion started at $0$. Since
$T_{a+b}$ is the sum of $T_a$ and of an independent copy of $T_b$,
\begin{equation}\label{eq:q-conv}
	\int_0^t q(a,s)\,q(b,t-s)\,ds = q(a+b,t), \qquad a,b>0 .
\end{equation}

\begin{lemma}\label{lem:sbm-abs}
	Let $X$ be a skew Brownian motion started at $x$, and $W'$ a Brownian motion
	started at $|x|$. Then the processes $\bigl(|X_t|,\hat{L}_t\bigr)_{t\geqslant0}$
	and $\bigl(|W'_t|,L^0_t(W')\bigr)_{t\geqslant0}$ have the same law.
\end{lemma}
\begin{proof}
	Apply the symmetric Itô--Tanaka formula (Theorem~\ref{th:symitotanaka}) to
	$f(z) = |z|$, whose symmetric derivative is $\sgn$ with $\sgn(0) = 0$ and
	with $f'' = 2\delta_0$:
	\[
		|X_t| = |x| + \int_0^t\sgn(X_s)\,dX_s + \hat{L}_t
		= |x| + \tilde{W}_t + \hat{L}_t,
		\qquad \tilde{W}_t = \int_0^t\sgn(X_s)\,dW_s .
	\]
	The term $\beta\int_0^t\sgn(X_s)\,d\hat{L}_s$ vanishes because $d\hat{L}$
	is carried by $\{s : X_s = 0\}$. Since $\{s : X_s = 0\}$ has Lebesgue measure
	$0$ by the occupation formula, $\langle\tilde{W}\rangle_t = t$, and
	$\tilde{W}$ is a Brownian motion by Lévy's characterisation. By Skorokhod's
	lemma (Lemma~\ref{lem:skorokhod}), applied to $y = |x| + \tilde{W}$,
	\[
		\hat{L}_t = \sup_{s\leqslant t}\bigl(-|x| - \tilde{W}_s\bigr)\vee0,
		\qquad |X_t| = |x| + \tilde{W}_t + \hat{L}_t .
	\]
	The same computation for $W'$ (for which $\hat{L}^0(W') = L^0(W')$) gives
	$|W'| = |x| + \tilde{W}' + L^0(W')$ with a Brownian motion $\tilde{W}'$, so
	both pairs are the same functional of a Brownian motion.
\end{proof}

\begin{lemma}\label{lem:bm-joint}
	Let $W$ be a Brownian motion started at $0$. For $t>0$,
	\[
		\mathbb{P}_0\bigl[|W_t|\in dy\,;\,L^0_t(W)\in dl\bigr] = 2\,q(y+l,t)\,dy\,dl,
		\qquad y>0,\ l>0 .
	\]
\end{lemma}
\begin{proof}
	Let $M_t = \sup_{s\leqslant t}W_s$. By Lévy's theorem, which follows from
	Skorokhod's lemma as in the proof of Lemma~\ref{lem:sbm-abs},
	$(|W_t|,L^0_t(W))$ has the law of $(M_t - W_t,M_t)$. By the reflection
	principle \citep[Section~2.8]{KaratzasShreve1991}, $(M_t,W_t)$ has density
	$\frac{2(2m-w)}{\sqrt{2\pi t^3}}\exp\bigl(-\frac{(2m-w)^2}{2t}\bigr)$ on
	$\{m>0,\ w<m\}$. The change of variables $l = m$, $y = m-w$ has Jacobian $1$
	and gives the result.
\end{proof}

The next lemma is intuitive, but the presence of the local time makes it hard
to prove directly from \eqref{eq:sbm}. The construction by random flipping
makes it easy.

\begin{lemma}\label{lem:sbm-flip}
	For $t>0$ and Borel sets $A\subset\,]0,\infty[$ and $C\subset[0,\infty[$,
	\begin{align*}
		\mathbb{P}_0\bigl[X_t\in A\,;\,\hat{L}_t\in C\bigr]
		 & = \frac{1+\beta}{2}\,\mathbb{P}_0\bigl[|X_t|\in A\,;\,\hat{L}_t\in C\bigr], \\
		\mathbb{P}_0\bigl[X_t\in -A\,;\,\hat{L}_t\in C\bigr]
		 & = \frac{1-\beta}{2}\,\mathbb{P}_0\bigl[|X_t|\in A\,;\,\hat{L}_t\in C\bigr].
	\end{align*}
\end{lemma}
\begin{proof}
	By uniqueness in law, and since $\hat{L}$ is a functional of the path of
	$X$, we may take $X$ to be the process constructed above. Since
	$\mathbb{P}(B_t = 0) = 0$, almost surely $t$ belongs to some excursion
	interval $I_{N}$, where the index $N$ is $\sigma(B)$-measurable. On
	$\{|B_t|\in A\}$, we have $X_t\in A$ if and only if $Y_N = 1$. Using
	$|X| = |B|$, $\hat{L}_t = L^0_t(B)$, and the independence of $(Y_n)$ and
	$B$,
	\begin{align*}
		\mathbb{P}\bigl[X_t\in A\,;\,\hat{L}_t\in C\bigr]
		 & = \E\Bigl[1_{\{|B_t|\in A,\ L^0_t(B)\in C\}}\,
		\mathbb{P}\bigl(Y_N = 1\,\big|\,\sigma(B)\bigr)\Bigr]                                   \\
		 & = \E\Bigl[1_{\{|B_t|\in A,\ L^0_t(B)\in C\}}\sum_{n}1_{\{N=n\}}\,\mathbb{P}(Y_n=1)\Bigr] \\
		 & = \frac{1+\beta}{2}\,\mathbb{P}\bigl[|X_t|\in A\,;\,\hat{L}_t\in C\bigr].
	\end{align*}
	The second identity is proved in the same way, with $Y_N = -1$.
\end{proof}

From these lemmas we can compute the joint law of $(X_t,\hat{L}_t)$
\citep{EtoreMartinez2012}.

\begin{proposition}\label{prop:sbm-joint}
	Let $x\geqslant0$ and $t>0$. Under $\mathbb{P}_x$, the law of
	$(X_t,\hat{L}_t)$ is
	\begin{align*}
		\mathbb{P}_x\bigl[X_t\in dy\,;\,\hat{L}_t\in dl\bigr]
		 & = 1_{\{y>0,\,l>0\}}\,(1+\beta)\,q(x+y+l,t)\,dy\,dl          \\
		 & \quad + 1_{\{y<0,\,l>0\}}\,(1-\beta)\,q(x-y+l,t)\,dy\,dl     \\
		 & \quad + 1_{\{y>0\}}\bigl(p_t(y-x) - p_t(y+x)\bigr)\,dy\,\delta_0(dl).
	\end{align*}
	For $x<0$, the law is obtained by applying this formula to $-X$, which is a
	skew Brownian motion of parameter $-\beta$ started at $-x$, with the same
	symmetric local time at $0$.
\end{proposition}
\begin{proof}
	\emph{Step 1: $x = 0$.} By Lemmas~\ref{lem:sbm-flip}, \ref{lem:sbm-abs}
	and~\ref{lem:bm-joint}, for $u>0$, $y>0$ and $l>0$,
	\[
		\mathbb{P}_0\bigl[X_u\in\pm dy\,;\,\hat{L}_u\in dl\bigr]
		= \frac{1\pm\beta}{2}\,\mathbb{P}_0\bigl[|W_u|\in dy\,;\,L^0_u(W)\in dl\bigr]
		= (1\pm\beta)\,q(y+l,u)\,dy\,dl .
	\]
	Moreover $\hat{L}_u>0$ for every $u>0$, $\mathbb{P}_0$-almost surely, since
	$L^0_u(W)$ has the law of $M_u>0$. This is the formula for $x=0$.

	\emph{Step 2: $x>0$, $l>0$.} Let $T_0 = \inf\{s : X_s = 0\}$. On
	$[0,T_0]$, $X = x + W$ and $\hat{L} = 0$, so the law of $T_0$ has density
	$q(x,\cdot)$. By the strong Markov property at $T_0$, and since
	$X_{T_0} = 0$ and $\hat{L}_{T_0+u} = \hat{L}^0_u(X_{T_0+\cdot})$, Step~1
	and~\eqref{eq:q-conv} give, for $y>0$,
	\begin{align*}
		\mathbb{P}_x\bigl[X_t\in dy\,;\,\hat{L}_t\in dl\bigr]
		 & = \int_0^t q(x,s)\,\mathbb{P}_0\bigl[X_{t-s}\in dy\,;\,\hat{L}_{t-s}\in dl\bigr]\,ds \\
		 & = (1+\beta)\int_0^t q(x,s)\,q(y+l,t-s)\,ds\;dy\,dl
		= (1+\beta)\,q(x+y+l,t)\,dy\,dl,
	\end{align*}
	and in the same way, for $y<0$, we get $(1-\beta)\,q(x+|y|+l,t)\,dy\,dl$.

	\emph{Step 3: $x>0$, $l=0$.} By Step~1, $\hat{L}_{T_0+u}>0$ for all $u>0$,
	so $\{\hat{L}_t = 0\} = \{t\leqslant T_0\}$, and
	$\mathbb{P}_x(T_0 = t) = 0$. On $\{t<T_0\}$, $X_t = x + W_t > 0$. By the
	reflection principle,
	\[
		\mathbb{P}_x\bigl[X_t\in dy\,;\,\hat{L}_t = 0\bigr]
		= \mathbb{P}_x\bigl[x+W_t\in dy\,;\,t<T_0\bigr]
		= 1_{\{y>0\}}\bigl(p_t(y-x) - p_t(y+x)\bigr)\,dy . \qedhere
	\]
\end{proof}

Integrating in $l$, using $\int_0^\infty q(a+l,t)\,dl = p_t(a)$, we obtain
the transition density of $X$.

\begin{proposition}\label{prop:sbm-density}
	For $t>0$, $x\in\mathbb{R}$ and $y\neq0$,
	\[
		\mathbb{P}_x[X_t\in dy] = p^\beta_t(x,y)\,dy, \qquad
		p^\beta_t(x,y) = p_t(y-x) + \beta\sgn(y)\,p_t(|x|+|y|).
	\]
\end{proposition}
\begin{proof}
	Let $x\geqslant0$. For $y>0$, Proposition~\ref{prop:sbm-joint} gives
	$(1+\beta)p_t(x+y) + p_t(y-x) - p_t(y+x) = p_t(y-x) + \beta p_t(x+y)$. For
	$y<0$, it gives $(1-\beta)p_t(x+|y|)$, and $x+|y| = |y-x|$, so this is
	$p_t(y-x) - \beta p_t(|x|+|y|)$. For $x<0$, apply the result to $-X$, which
	replaces $(x,y,\beta)$ by $(-x,-y,-\beta)$ and leaves the formula unchanged.
\end{proof}

For $\beta = 0$ this is the Brownian transition density. Since
$\int\sgn(y)\,p_t(|x|+|y|)\,dy = 0$, $p^\beta_t(x,\cdot)$ has total mass
$1$. Taking $x = 0$ gives $\mathbb{P}_0(X_t>0) = \frac{1+\beta}{2}$, as the
construction by random flipping suggests.

\section{The partial differential equation solved by skew Brownian motion}\label{sec:sbm-pde}

The drift $\beta\hat{L}^0(X)$ in \eqref{eq:sbm} is singular, so the generator
of $X$ is not a classical differential operator. Formally it is
\[
	\mathcal{A} = \frac{1}{2}\frac{d^2}{dx^2} + \beta\,\delta_0\,\frac{d}{dx}.
\]
\citet{Portenko1990} identified processes of this type through such generalised
operators, and developed their links with the partial differential equation
that the transition semigroup solves. This link gives several methods to
compute the semigroup of skew Brownian motion. As in Chapter~\ref{ch:pde}, the
probabilistic solution can be checked directly.

\begin{proposition}\label{prop:sbm-pde}
	Let $f$ be bounded and continuous, and let $u(t,x) = \E_x f(X_t)$. Then
	\begin{enumerate}[label=(\roman*)]
		\item $u\in C^{1,2}\bigl(]0,\infty[\times(\mathbb{R}\setminus\{0\})\bigr)$ and
		      $\dfrac{\partial u}{\partial t} = \dfrac{1}{2}\dfrac{\partial^2u}{\partial x^2}$
		      on $]0,\infty[\times(\mathbb{R}\setminus\{0\})$;
		\item $u$ is continuous on $[0,\infty[\times\mathbb{R}$, and $u(0,x) = f(x)$;
		\item for $t>0$, $\frac{\partial u}{\partial x}(t,\cdot)$ has limits at
		      $0^+$ and $0^-$, and
		      \[
			      (1+\beta)\frac{\partial u}{\partial x}(t,0^+) = (1-\beta)\frac{\partial u}{\partial x}(t,0^-).
		      \]
	\end{enumerate}
\end{proposition}
\begin{proof}
	By Proposition~\ref{prop:sbm-density},
	\[
		u(t,x) = \int p_t(y-x)\,f(y)\,dy + \beta\int\sgn(y)\,p_t(|x|+|y|)\,f(y)\,dy .
	\]
	On each half-line $x>0$ and $x<0$, both kernels $(t,x)\mapsto p_t(y-x)$ and
	$(t,x)\mapsto p_t(|x|+|y|) = p_t(\pm x + |y|)$ are solutions of the heat
	equation. Differentiation under the integral sign is justified exactly as in
	Theorem~\ref{th415}, which gives (i). The same domination shows that the
	$x$-derivatives extend continuously to $x = 0^\pm$. Using
	$p_t'(z) = -\frac{z}{t}p_t(z)$, we get
	\[
		\frac{\partial u}{\partial x}(t,0^\pm)
		= \int\frac{y}{t}\,p_t(y)\,f(y)\,dy
		\mp\beta\int\sgn(y)\,\frac{|y|}{t}\,p_t(y)\,f(y)\,dy
		= (1\mp\beta)\,a_t,
	\]
	where $a_t = \int\frac{y}{t}\,p_t(y)\,f(y)\,dy$. Hence
	$(1+\beta)\partial_xu(t,0^+) = (1+\beta)(1-\beta)a_t = (1-\beta)\partial_xu(t,0^-)$,
	which is (iii).

	For (ii), continuity on $]0,\infty[\times\mathbb{R}$ is clear from the
	formula. As $t\to0$ and $x\to x_0$, the first term tends to $f(x_0)$ by
	Theorem~\ref{th414}. The second term is bounded by
	$\|f\|_\infty\int p_t(|x|+|y|)\,dy$, which tends to $0$ if $x_0\neq0$. If
	$x_0 = 0$, it equals
	$\beta\int_0^\infty p_t(|x|+z)\bigl(f(z)-f(-z)\bigr)\,dz$, which tends to
	$0$ because $f(z)-f(-z)\to0$ as $z\to0$, while the measures
	$p_t(|x|+z)\,dz$ on $]0,\infty[$ have mass at most $\frac12$ and concentrate
	near $z = 0$.
\end{proof}

Thus $u$ solves the heat equation on each side of $0$, is continuous across
$0$, and satisfies the \emph{transmission condition}
\[
	(1+\beta)\frac{\partial u}{\partial x}(t,0^+) = (1-\beta)\frac{\partial u}{\partial x}(t,0^-).
\]
This condition at $0$ involves the skewness parameter, and it makes the
equation harder to solve than the heat equation. We now explain where it comes
from, and then give an equivalent formulation in divergence form.

\subsection{The transmission condition}
Let $\mathcal{D}$ be the set of continuous functions $f$ that are $C^2$ on
$\mathbb{R}\setminus\{0\}$, with $f'$ and $f''$ bounded and having limits at
$0^+$ and $0^-$. Such an $f$ is a difference of two convex functions, and in
the sense of distributions
\[
	f''(dz) = f''(z)\,dz + \bigl(f'(0^+) - f'(0^-)\bigr)\,\delta_0(dz),
\]
while its symmetric derivative at $0$ is $\frac{1}{2}\bigl(f'(0^+)+f'(0^-)\bigr)$.
By the symmetric Itô--Tanaka formula (Theorem~\ref{th:symitotanaka}),
\[
	f(X_t) - f(x) = \int_0^t f'_{\mathrm{s}}(X_s)\,dW_s
	+ \beta\int_0^t f'_{\mathrm{s}}(X_s)\,d\hat{L}_s
	+ \frac{1}{2}\int_{\mathbb{R}}\hat{L}^a_t\,f''(da) .
\]
Since $d\hat{L}$ is carried by $\{X = 0\}$, the second term is
$\frac{\beta}{2}\bigl(f'(0^+)+f'(0^-)\bigr)\hat{L}_t$. By the occupation
formula, the last term is
$\frac{1}{2}\int_0^tf''(X_s)\,ds + \frac{1}{2}\bigl(f'(0^+)-f'(0^-)\bigr)\hat{L}_t$.
Hence
\begin{equation}\label{eq:sbm-ito}
	f(X_t) - f(x) - \frac{1}{2}\int_0^tf''(X_s)\,ds
	= \int_0^tf'_{\mathrm{s}}(X_s)\,dW_s
	+ \frac{1}{2}\Bigl((1+\beta)f'(0^+) - (1-\beta)f'(0^-)\Bigr)\hat{L}_t .
\end{equation}
The stochastic integral is a martingale, because $f'$ is bounded. The last
term is a continuous process of bounded variation, and $\hat{L}$ is not
identically $0$ since $X$ reaches $0$. Hence the left-hand side of
\eqref{eq:sbm-ito} is a martingale if and only if
\[
	(1+\beta)f'(0^+) = (1-\beta)f'(0^-).
\]
In this sense the generator of $X$ acts as $f\mapsto\frac{1}{2}f''$ on the
functions of $\mathcal{D}$ that satisfy the transmission condition; see
\citet{Lejay2006} for a precise description of its domain. For $f = F_\nu$,
the scale function of Remark~\ref{rem:sbm-scale}, $f'' = 0$ on
$\mathbb{R}\setminus\{0\}$, and \eqref{eq:sbm-ito} reduces to the martingale
property of $F_\nu(X)$.

\subsection{Divergence form}
Assume $|\beta|<1$, and let
\[
	k(x) = \frac{1-\beta}{2}\,1_{\{x<0\}} + \frac{1+\beta}{2}\,1_{\{x>0\}},
	\qquad
	L = \frac{1}{2k(x)}\frac{d}{dx}\Bigl(k(x)\frac{d}{dx}\Bigr),
\]
with domain
$\mathcal{D}(L) = \{\psi\in H^1(\mathbb{R}) : k\psi'\in H^1(\mathbb{R})\}$.
A function $\psi\in\mathcal{D}(L)$ is continuous, and so is $k\psi'$; at $0$
this is exactly $(1+\beta)\psi'(0^+) = (1-\beta)\psi'(0^-)$. Away from $0$,
$L\psi = \frac{1}{2}\psi''$. So $L$ is the generator of skew Brownian motion,
written in divergence form \citep{Lejay2006}.

The operator $L$ is self-adjoint and nonpositive on $L^2(\mathbb{R},k\,dx)$:
for $\psi,\varphi\in\mathcal{D}(L)$,
\[
	\int L\psi\,\varphi\,k\,dx = -\frac{1}{2}\int k\,\psi'\varphi'\,dx .
\]
It is associated with the Dirichlet form
$\mathcal{E}(\psi,\varphi) = \frac{1}{2}\int k\,\psi'\varphi'\,dx$ on
$H^1(\mathbb{R})$. Since $k$ is bounded above and below by positive constants,
$L^2(\mathbb{R},k\,dx)$ and $L^2(\mathbb{R},dx)$ coincide as sets and have
equivalent norms, but $L$ is symmetric only for the weighted scalar product.
The transmission problem can then be replaced by its weak formulation: for
$\phi\in L^2(\mathbb{R},k\,dx)$, find
\[
	u\in C\bigl(\mathbb{R}_+;L^2(\mathbb{R},k\,dx)\bigr)\cap L^2\bigl(\mathbb{R}_+;H^1(\mathbb{R})\bigr),
	\qquad \frac{\partial u}{\partial t} = Lu,\qquad u(0,\cdot) = \phi .
\]
The solution is $u(t,\cdot) = e^{tL}\phi$. An advantage of the divergence form
is that the Markov property of the associated process follows directly from
the theory of Dirichlet forms.
